\documentclass[10pt,a4paper]{amsart}
\usepackage[margin=19mm]{geometry}
\usepackage{amsmath,amssymb,mathtools}
\usepackage[scr=rsfs,cal=euler]{mathalfa}
\usepackage{xfrac}
\usepackage[unicode,hidelinks,bookmarks=false]{hyperref}
\newcommand{\ie}{i.e.\ }
\newcommand{\cf}{cf.\ }
\newcommand{\resp}{respectively\ }
\newcommand{\Subsection}{\subsection}
\providecommand{\nobreakdash}{\nobreak\hbox{-}}
\setlength{\emergencystretch}{2em}
\multlinegap0pt
\usepackage{amssymb}
\newtheorem{theorem}{Theorem}
\title{Haar-Bayesian Pure-State Prediction under Relative-Entropy Loss: Arbitrary-Effect Reduction and Global Optimality}
\author{Masahito Hayashi\textsuperscript{$*$}, Ayanava Dasgupta\textsuperscript{$\dagger$} and Naqueeb Ahmad Warsi\textsuperscript{$\dagger$}}
\date{Source: arXiv:2609.10072v1 (CC BY 4.0)}
\begin{document}
\maketitle

\noindent\emph{Source and licence.} Haar-Bayesian Pure-State Prediction under Relative-Entropy Loss: Arbitrary-Effect Reduction and Global Optimality. Version arXiv:2609.10072v1 (CC BY 4.0), distributed by arXiv under the Creative Commons Attribution 4.0 International licence, http://creativecommons.org/licenses/by/4.0/, as recorded in the arXiv licence metadata for this exact version and shown on the abstract page https://arxiv.org/abs/2609.10072v1. Full licence notice: \texttt{licenses/arxiv-2609.10072-CC-BY-4.0.txt}.\par\smallskip

\noindent\textit{Editorial scope.} Counted unit: the article's theorem thm:covariant\_reduction (source0.tex lines 747--774). The article's complete original proof of it is delivered with it (source0.tex lines 776--896, one \texttt{proof} environment of 121 lines). No mathematics is written, repaired or re-derived: the statement and the proof are the article's own lines, in the article's own notation, and no retained line is substituted or re-flowed.  No further source-local context is needed: every cross-reference of every retained line resolves inside the two delivered spans.  Nothing was omitted from any retained span, so the package carries no omission record.  No retained line contains a citation, so the wrapper carries no bibliography and no bibliography entry.  No retained line contains a figure environment, an image-inclusion command or drawing code, so no image is delivered and the package declares no asset dependency.  Editorial formatting only, in addition: an AMS \texttt{amsart} preamble that restates the article's own declarations and macros used by the retained lines, namely the environments \texttt{theorem} on separate counters, together with the standard packages those lines need.  The retained environments print in this document as Theorem 1 in order of appearance, whereas the article prints the same items with the numbering of its own full document.  The complete native source of the article as distributed in the arXiv e-print for this exact version is bundled unchanged as \texttt{sources/209872/source0.tex}.
\begin{theorem}[Hunt--Stein reduction for the present Bayes problem]
\label{thm:covariant_reduction}
For the Haar-prior pure-state model with relative-entropy loss, the minimum
Haar-Bayes risk over all decision procedures is equal to the minimum Haar-Bayes
risk over covariant decision procedures.  Moreover, for every decision POVM
\(N\), there exists a covariant decision POVM \(\overline N\) such that
\[
    R_{\mu_\Theta}(\overline N)=R_{\mu_\Theta}(N).
\]
For a covariant decision POVM \(\overline N\), the risk is constant on
\(\Theta\), and hence
\[
    R_{\max}(\overline N)=R_{\mu_\Theta}(\overline N).
\]
In particular, if a covariant procedure \(N_\star\) attains the Bayes optimum
\(B\), then its constant frequentist risk equals \(B\), and
\[
    B
    \le
    \inf_N R_{\max}(N)
    \le
    R_{\max}(N_\star)
    =
    R_{\mu_\Theta}(N_\star)
    =B.
\]
Hence \(N_\star\) is also minimax and the minimax value equals \(B\).
\end{theorem}

\begin{proof}
Fix an
arbitrary decision POVM \(N\).  For each \(g\in G\), define the rotated decision
POVM
\[
    N_g(B)
    :=
    U_g^{\otimes n\dagger}N(gB)U_g^{\otimes n},
    \qquad B\subset\mathcal A.
    \label{eq:rotated_decision_povm_def}
\]
We first compute its risk.  Using the change of decision variable
\(\omega=g\sigma\), the covariance of the state family, and the
loss-invariance relation displayed above, we obtain
\[
\begin{aligned}
    R_\theta(N_g)
    &=
    \int_{\mathcal A}
    \ell(\theta,\sigma)
    \operatorname{Tr}\!\left[
        \rho_\theta^{\otimes n}
        U_g^{\otimes n\dagger}N(gd\sigma)U_g^{\otimes n}
    \right]                                                   \\
    &=
    \int_{\mathcal A}
    \ell(\theta,g^{-1}\omega)
    \operatorname{Tr}\!\left[
        (U_g^{\otimes n}\rho_\theta^{\otimes n}U_g^{\otimes n\dagger})
        N(d\omega)
    \right]                                                   \\
    &=
    \int_{\mathcal A}
    \ell(g\theta,\omega)
    \operatorname{Tr}\!\left[
        \rho_{g\theta}^{\otimes n}N(d\omega)
    \right]                                                   
    =
    R_{g\theta}(N).
\end{aligned}
\]
Averaging this identity over \(\theta\) and using invariance of
\(\mu_\Theta\), we get
\[
    R_{\mu_\Theta}(N_g)
    =
    \int_\Theta R_{g\theta}(N)\,\mu_\Theta(d\theta)
    =
    R_{\mu_\Theta}(N).
\]

Now define the averaged POVM
\[
    \overline N(B)
    :=
    \int_G N_g(B)\,\mu_G(dg),
    \label{eq:averaged_decision_povm_def}
\]
where \(\mu_G\) is the normalized Haar measure on the compact group \(G\).  The
operator integral is well defined in the weak sense; positivity and
normalization follow from positivity and normalization of each \(N_g\).
Linearity of the risk in the POVM gives
\[
\begin{aligned}
    R_{\mu_\Theta}(\overline N)
    &=
    \int_G R_{\mu_\Theta}(N_g)\,\mu_G(dg) 
    =
    R_{\mu_\Theta}(N),
\end{aligned}
\]
where the last equality is the preceding equality.  Hence averaging
does not change the Haar-Bayes risk.

It remains to verify covariance.  For \(h\in G\),
\[
\begin{aligned}
    \overline N(hB)
    &=
    \int_G
    U_g^{\otimes n\dagger}N(ghB)U_g^{\otimes n}
    \,\mu_G(dg).
\end{aligned}
\]
Set \(k=gh\).  By right invariance of Haar measure, \(\mu_G(dg)=\mu_G(dk)\), and
\[
\begin{aligned}
    \overline N(hB)
    &=
    \int_G
    U_{kh^{-1}}^{\otimes n\dagger}N(kB)U_{kh^{-1}}^{\otimes n}
    \,\mu_G(dk)                                                \\
    &=
    U_h^{\otimes n}
    \left(
        \int_G U_k^{\otimes n\dagger}N(kB)U_k^{\otimes n}
        \,\mu_G(dk)
    \right)
    U_h^{\otimes n\dagger}                                      \\
    &=
    U_h^{\otimes n}\overline N(B)U_h^{\otimes n\dagger}.
\end{aligned}
\]
Thus \(\overline N\) is covariant in the sense of
the covariance condition displayed above.

Finally, suppose \(N\) is covariant.  Then the first calculation above, applied
to \(N\) itself, gives
\[
    R_{g\theta}(N)=R_\theta(N)
\]
for all \(g\in G\).  Since \(G\) acts transitively on \(\Theta\), the risk is
constant on \(\Theta\).  Hence
\[
    R_{\max}(N)=R_{\mu_\Theta}(N).
\]
Combining this with the averaging construction proves that the Haar-Bayes
optimization may be restricted to covariant decision POVMs, and that for
covariant procedures the Haar-Bayes and worst-case risks coincide.

\end{proof}

\end{document}
